\subsection{Aplicación a los morfismos propios: I. La propiedad de cambio de base}
\label{subsection:III.7.7}
\label{III.7.7}

Los tres números que siguen son, en lo esencial, traducciones al lenguaje de
los morfismos de preesquemas de los resultados de los números precedentes.

\begin{env}[7.7.1]
\label{III.7.7.1}
Sea $f:X\to Y$ un morfismo cuasicompacto y separado de preesquemas, y sea
$\mathscr P_\bullet$ un complejo de $\mathscr O_X$-módulos cuasicoherentes cuyo
operador de derivación es de grado $-1$; supongamos además que los
$\mathscr O_X$-módulos $\mathscr P_i$ son planos sobre $Y$
\sref[0]{0.6.7.1}.

\oldpage[III]{198}
Consideraremos el $\partial$-funtor
$\mathscr M\mapsto\mathscr T_\bullet(\mathscr M)$ (también denotado por
$\mathscr T_\bullet(\mathscr P_\bullet,\mathscr M)$) de la categoría de los
$\mathscr O_Y$-módulos cuasicoherentes en la categoría de los
$\mathscr O_Y$-módulos cuasicoherentes, en virtud de \sref{III.6.2.3}, definido
por
\[
  \mathscr T_n(\mathscr P_\bullet,\mathscr M)
  =\mathscr T_n(\mathscr M)
  =\mathscr H^{-n}
    \bigl(f,\mathscr P^\bullet\otimes_{\mathscr O_Y}\mathscr M\bigr)
  \qquad\text{para }n\in\mathbf Z,
  \tag{7.7.1.1}\label{III.7.7.1.1}
\]
donde $\mathscr P^\bullet$ es el complejo cuyo término de grado $j$ es
$\mathscr P_{-j}$, de modo que su operador de derivación es entonces de grado
$+1$. El funtor $\mathscr T_\bullet$ así definido es un \emph{funtor
homológico} en $\mathscr M$ \sref{III.6.2.6}.
\end{env}

\begin{env}[7.7.2]
\label{III.7.7.2}
Sea $g:Y'\to Y$ un morfismo, y pongamos
\[
  X'=X_{(Y')}=X\times_Y Y',
  \qquad
  f'=f_{(Y')}:X'\longrightarrow Y',
\]
que es un morfismo cuasicompacto y separado; pongamos además
\[
  \mathscr P'_\bullet
  =\mathscr P_\bullet\otimes_{\mathscr O_Y}\mathscr O_{Y'}.
\]
Este es un complejo de $\mathscr O_{X'}$-módulos cuasicoherentes que son planos
sobre $Y'$, en virtud de \sref[I]{I.9.1.12} y \sref[0]{0.6.2.1}. Pondremos,
con las mismas convenciones sobre los grados,
\[
  \mathscr T^{Y'}_\bullet(\mathscr M')
  =\mathscr H^\bullet
    \bigl(f',\mathscr P'^\bullet\otimes_{\mathscr O_{Y'}}\mathscr M'\bigr)
  =\mathscr H^\bullet
    \bigl(f',\mathscr P^\bullet\otimes_{\mathscr O_Y}\mathscr M'\bigr)
  \tag{7.7.2.1}\label{III.7.7.2.1}
\]
que es un funtor homológico en el $\mathscr O_{Y'}$-módulo cuasicoherente
$\mathscr M'$. Cuando $Y'$ sea un esquema afín de anillo $A'$, escribiremos
$\mathscr T^{A'}_\bullet$ en lugar de $\mathscr T^{Y'}_\bullet$; para todo
$A'$-módulo $M'$ se tiene entonces
\[
  \mathscr T^{A'}_\bullet(\widetilde{M'})
  =\bigl(\Gamma(Y',\mathscr T^{Y'}_\bullet(\widetilde{M'}))\bigr)^\sim.
\]
Pondremos
\[
  T^{A'}_\bullet(M')
  =\Gamma(Y',\mathscr T^{Y'}_\bullet(\widetilde{M'})),
\]
que es un funtor homológico de $A'$-módulos con valores en la categoría de los
$A'$-módulos.

Obsérvese que, si $Y=\operatorname{Spec}(A)$ es también afín, el funtor de
$A'$-módulos $T^{A'}_\bullet$ coincide con el funtor obtenido por extensión de
escalares de $A$ a $A'$ a partir del funtor homológico de $A$-módulos
$T^A_\bullet$ \sref{III.7.1.3}. En efecto, sea $g:Y'\to Y$ el morfismo
correspondiente al homomorfismo de anillos $A\to A'$, y sea $g':X'\to X$ el
morfismo correspondiente, que es afín \sref[II]{II.1.6.2}; si $\mathfrak U$ es
un recubrimiento abierto afín de $X$, entonces
$\mathfrak U'=g'^{-1}(\mathfrak U)$ es un recubrimiento abierto afín de $X'$.
En virtud de \sref{III.6.2.2}, todo se reduce a ver que
\[
  C^\bullet
    \bigl(\mathfrak U,
      \mathscr P_\bullet\otimes_{\mathscr O_Y}g_*(\mathscr M')\bigr)
  =C^\bullet
    \bigl(\mathfrak U',
      \mathscr P_\bullet\otimes_{\mathscr O_Y}\mathscr M'\bigr),
\]
y, finalmente, que para todo abierto afín $U$ de $X$, poniendo
$U'=g'^{-1}(U)$, se tiene
\[
  \Gamma
    \bigl(U,\mathscr P_\bullet\otimes_{\mathscr O_Y}g_*(\mathscr M')\bigr)
  =\Gamma
    \bigl(U',\mathscr P_\bullet\otimes_{\mathscr O_Y}\mathscr M'\bigr),
\]
lo cual es trivial \sref[I]{I.1.3}, \sref[I]{I.3.2}.

En particular, si $U$ es un abierto de $Y$, para todo
$\mathscr O_Y$-módulo cuasicoherente $\mathscr M$ se tiene
\[
  \mathscr T^U_\bullet(\mathscr M|U)
  =\bigl(\mathscr T_\bullet(\mathscr M)\bigr)|U.
  \tag{7.7.2.2}\label{III.7.7.2.2}
\]
\end{env}

\begin{env}[7.7.3]
\label{III.7.7.3}
Para todo $\mathscr O_Y$-módulo cuasicoherente $\mathscr M$ existe un
homomorfismo canónico, funtorial en $\mathscr M$,
\[
  \mathscr T_p(\mathscr O_Y)\otimes_{\mathscr O_Y}\mathscr M
  \longrightarrow
  \mathscr T_p(\mathscr M).
  \tag{7.7.3.1}\label{III.7.7.3.1}
\]
En efecto, si $Y$ es afín, este homomorfismo se definió en
\sref{III.7.2.2}; esta definición se extiende sin dificultad al caso general,
observando que, si $U$ y $V$ son dos abiertos afines de $Y$ tales que
$V\subset U$, el diagrama

\oldpage[III]{199}
\[
\xymatrix@C=18pt@R=32pt{
  {\begin{gathered}
    \bigl(\mathscr T_p(\mathscr O_Y)
      \otimes_{\mathscr O_Y}\mathscr M\bigr)|U\\
    =\mathscr T^U_p(\mathscr O_Y|U)
      \otimes_{\mathscr O_Y|U}(\mathscr M|U)
  \end{gathered}}
  \ar[r]\ar[d]
  &
  {\begin{gathered}
    \mathscr T^U_p(\mathscr M|U)\\
    =\bigl(\mathscr T_p(\mathscr M)\bigr)|U
  \end{gathered}}
  \ar[d]
  \\
  {\begin{gathered}
    \bigl(\mathscr T_p(\mathscr O_Y)
      \otimes_{\mathscr O_Y}\mathscr M\bigr)|V\\
    =\mathscr T^V_p(\mathscr O_Y|V)
      \otimes_{\mathscr O_Y|V}(\mathscr M|V)
  \end{gathered}}
  \ar[r]
  &
  {\begin{gathered}
    \mathscr T^V_p(\mathscr M|V)\\
    =\bigl(\mathscr T_p(\mathscr M)\bigr)|V
  \end{gathered}}
}
\]
es conmutativo por \sref{III.7.2.3.3}.

Para todo morfismo $g:Y'\to Y$ existe un homomorfismo canónico
\[
  \mathscr T_p(\mathscr O_Y)\otimes_{\mathscr O_Y}\mathscr O_{Y'}
  \longrightarrow
  \mathscr T^{Y'}_p(\mathscr O_{Y'}),
  \tag{7.7.3.2}\label{III.7.7.3.2}
\]
que no es sino el caso particular de \sref{III.6.7.11.2}, para los objetos
límite, en el caso en que $S=Y$, $v_1=f$, $v_2=1_Y$ y
$\mathscr P^{(2)}_\bullet$ se reduce al único término $\mathscr M$ de grado
$0$.

Cuando $Y=\operatorname{Spec}(A)$ e $Y'=\operatorname{Spec}(A')$ son afines,
\eqref{III.7.7.3.2} no es sino el homomorfismo de haces correspondiente al
homomorfismo canónico de $A'$-módulos definido en \sref{III.7.2.2},
\[
  T^A_p(A)\otimes_A A'
  \longrightarrow
  T^{A'}_p(A')=T^A_p(A'),
\]
como resulta fácilmente de \sref{III.6.7.11}, pues en el caso considerado se
puede tomar $\mathfrak U'^{(i)}=u_i^{-1}(\mathfrak U^{(i)})$.
\end{env}

\begin{env}[7.7.4]
\label{III.7.7.4}
Cuando $f$ es un morfismo \emph{propio}, $Y=\operatorname{Spec}(A)$ es un
esquema afín noetheriano y $\mathscr P_\bullet$ es un complejo acotado
inferiormente de $\mathscr O_X$-módulos coherentes y planos sobre $Y$, vimos en
\sref{III.6.10.5} que se puede escribir, salvo isomorfismo,
\[
  \mathscr T_p(\mathscr M)
  =\mathscr H_p
    \bigl(\mathscr L_\bullet\otimes_{\mathscr O_Y}\mathscr M\bigr),
  \qquad
  \mathscr L_\bullet=\widetilde{L_\bullet},
\]
donde $L_\bullet$ es un complejo acotado inferiormente de $A$-módulos libres
de tipo finito; por tanto, el funtor $\mathscr T_\bullet$ es del tipo estudiado
en detalle en \sref{III.7.4} y \sref{III.7.6}. Traduciremos ahora los resultados
de este estudio.
\end{env}

\begin{theorem}[7.7.5]
\label{III.7.7.5}
Sean $Y$ un preesquema localmente noetheriano, $(U_\alpha)$ un recubrimiento
de $Y$ formado por abiertos afines, $f:X\to Y$ un morfismo propio y
$\mathscr P_\bullet$ un complejo acotado inferiormente de
$\mathscr O_X$-módulos coherentes y planos sobre $Y$. El funtor homológico
$\mathscr T_\bullet(\mathscr M)$ definido por \eqref{III.7.7.1.1} posee
entonces las propiedades siguientes:

\medskip
\noindent
\textnormal{I) (La propiedad de semicontinuidad).}\footnote{Un caso particular
de este teorema se encuentra ya en la nota [3] de Chow--Igusa. La propiedad de
semicontinuidad fue descubierta, en el marco de los espacios analíticos y bajo
hipótesis bastante particulares, por Kodaira--Spencer (``On the variations of
almost-complex structures'', \emph{Algebraic Geometry and Topology, A Symposium
in honor of S. Lefschetz}, Princeton Series n\textsuperscript{o} 12,
p.~139--150, Princeton, 1957), y la versión general fue demostrada por Grauert
[5].}
\emph{La función}
\[
  y\longmapsto
  d_p(y)
  =
  \operatorname{rank}_{\kres(y)}
    \mathscr T_p^{\kres(y)}(\kres(y))
  \tag{7.7.5.1}\label{III.7.7.5.1}
\]
\emph{es semicontinua superiormente.}

\oldpage[III]{200}
\medskip
\noindent
\textnormal{II) (La propiedad de cambio de base).}
\emph{Para un entero $p$ dado, las condiciones siguientes son equivalentes:}
\begin{enumerate}
  \item[a)] $\mathscr T_p$ es exacto por la derecha.

  \item[a$'$)] $\mathscr T_p(\mathscr M)$ es isomorfo a un funtor de la forma
    $\mathscr N\otimes_{\mathscr O_Y}\mathscr M$, siendo $\mathscr N$
    necesariamente isomorfo a
    \[
      \mathscr T_p(\mathscr O_Y)
      =
      \mathscr H^{-p}(f,\mathscr P^\bullet).
    \]

  \item[a$''$)] El homomorfismo canónico funtorial
    \eqref{III.7.7.3.1} es un isomorfismo.

  \item[b)] $\mathscr T_{p-1}$ es exacto por la izquierda.

  \item[b$'$)] Existe un $\mathscr O_Y$-módulo $\mathscr Q$ (necesariamente
    coherente y determinado salvo isomorfismo único) y un isomorfismo de
    funtores
    \[
      \mathscr T_{p-1}(\mathscr M)
      \xrightarrow{\sim}
      \mathscr{H}\!om_{\mathscr O_Y}(\mathscr Q,\mathscr M).
      \tag{7.7.5.2}\label{III.7.7.5.2}
    \]

  \item[c)] Denotando por $A_\alpha$ el anillo del abierto afín $U_\alpha$,
    para todo índice $\alpha$ el funtor de $A_\alpha$-módulos
    $T_p^{A_\alpha}$ es exacto por la derecha.

  \item[d)] Para todo morfismo $g:Y'\to Y$, el homomorfismo canónico
    \[
      \mathscr T_p(\mathscr O_Y)
      \otimes_{\mathscr O_Y}\mathscr O_{Y'}
      \longrightarrow
      \mathscr T_p^{Y'}(\mathscr O_{Y'})
      \tag{7.7.5.3}\label{III.7.7.5.3}
    \]
    es un isomorfismo.
\end{enumerate}
\end{theorem}

\begin{proof}
La propiedad de semicontinuidad es local sobre $Y$ y resulta, por tanto, de
\sref{III.7.7.4} y \sref{III.7.6.9}. Es claro que a$''$) implica a$'$) y que
a$'$) implica a). La equivalencia de a), a$''$), b) y b$'$) se demostró en
\sref{III.7.3.1} y \sref{III.7.4.6}, habida cuenta de
\sref{III.7.7.4}, cuando $Y$ es afín.

Para pasar al caso general, probemos primero que a) equivale a c), lo cual
probará que la propiedad a) es local sobre $Y$; la demostración se aplica
igualmente para probar el carácter local de a$''$) y b). Como es claro que c)
implica a), todo se reduce a probar el recíproco. Basta evidentemente demostrar
que, para todo abierto afín $U$ de $Y$ y toda sucesión exacta
\[
  0\longrightarrow\mathscr F'
  \longrightarrow\mathscr F
  \longrightarrow\mathscr F''
  \longrightarrow0
\]
de $(\mathscr O_Y|U)$-módulos cuasicoherentes, existe una sucesión exacta
\[
  0\longrightarrow\mathscr G'
  \longrightarrow\mathscr G
  \longrightarrow\mathscr G''
  \longrightarrow0
\]
de $\mathscr O_Y$-módulos cuasicoherentes tal que
\[
  \mathscr F'=\mathscr G'|U,
  \qquad
  \mathscr F=\mathscr G|U,
  \qquad
  \mathscr F''=\mathscr G''|U.
\]
Ahora bien, esto resulta inmediatamente de la hipótesis de que $Y$ es
localmente noetheriano y de \sref[I]{I.9.4.2}: en efecto, se prolonga
$\mathscr F$ a un $\mathscr O_Y$-módulo cuasicoherente $\mathscr G$,
$\mathscr F'$ a un sub-$\mathscr O_Y$-módulo $\mathscr G'$ de $\mathscr G$, y
basta tomar $\mathscr G''=\mathscr G/\mathscr G'$.

Para demostrar la equivalencia de b) y b$'$) en el caso general, observemos que,
cuando $Y$ es afín, sabemos que $\mathscr Q$ está determinado salvo isomorfismo
único; si entonces $U$ es un abierto afín del esquema afín $Y$, se deduce que
existe un isomorfismo funtorial
\[
  \mathscr T_{p-1}^{U}(\mathscr M|U)
  \xrightarrow{\sim}
  \mathscr{H}\!om_{\mathscr O_Y|U}
    (\mathscr Q|U,\mathscr M|U).
\]
En el caso general, para todo abierto afín $U$ de $Y$ existe un
$(\mathscr O_Y|U)$-módulo coherente $\mathscr Q_U$ y un isomorfismo funtorial
\[
  \mathscr T_{p-1}^{U}(\mathscr M|U)
  \xrightarrow{\sim}
  \mathscr{H}\!om_{\mathscr O_Y|U}
    (\mathscr Q_U,\mathscr M|U).
\]
La observación precedente muestra que, si $V$ es un abierto afín contenido en
$U$, se tiene $\mathscr Q_U|V=\mathscr Q_V$; de ahí la existencia y unicidad
del $\mathscr O_Y$-módulo $\mathscr Q$ que verifica \eqref{III.7.7.5.2}.

Queda por demostrar la equivalencia de a) y d). Es claro que d) es local sobre
$Y$, y acabamos de ver que lo mismo ocurre con a); además, d) es también local
sobre $Y'$. Ahora bien, cuando $Y=\operatorname{Spec}(A)$ e
$Y'=\operatorname{Spec}(A')$, vimos que $T_\bullet^{A'}$ es el funtor obtenido
\oldpage[III]{201}
a partir de $T_\bullet^A$ por extensión de escalares a $A'$, y es claro entonces
que a$'$) implica que \eqref{III.7.7.5.3} es un isomorfismo.
Recíprocamente, supongamos todavía que $Y=\operatorname{Spec}(A)$ es afín y sea
$A'$ la $A$-álgebra $A\oplus M$, donde $M$ es un $A$-módulo cualquiera y la
multiplicación en $A'$ viene dada por
\[
  (a_1,m_1)(a_2,m_2)
  =
  (a_1a_2,a_1m_2+a_2m_1).
\]
Entonces
\[
  T_p^{A'}(A')
  =
  T_p(A\oplus M)
  =
  T_p(A)\oplus T_p(M).
\]
La hipótesis de que \eqref{III.7.7.5.3} sea biyectivo implica que lo mismo
ocurre con la aplicación canónica
\[
  T_p(A)\otimes_A M\longrightarrow T_p(M).
\]
En otras palabras, d) implica a$''$), lo que termina la demostración.
\end{proof}

\begin{theorem}[7.7.6]
\label{III.7.7.6}
Sean $Y$ un preesquema localmente noetheriano, $f:X\to Y$ un morfismo propio y
$\mathscr F$ un $\mathscr O_X$-módulo coherente y plano sobre $Y$. Existe
entonces un $\mathscr O_Y$-módulo coherente $\mathscr Q$ (determinado salvo
isomorfismo único) y un isomorfismo de funtores en el
$\mathscr O_Y$-módulo cuasicoherente $\mathscr M$,
\[
  f_*(\mathscr F\otimes_{\mathscr O_Y}\mathscr M)
  \xrightarrow{\sim}
  \mathscr{H}\!om_{\mathscr O_Y}(\mathscr Q,\mathscr M).
  \tag{7.7.6.1}\label{III.7.7.6.1}
\]
De ello se deduce un isomorfismo de funtores
\[
  \Gamma(X,\mathscr F\otimes_{\mathscr O_Y}\mathscr M)
  \xrightarrow{\sim}
  \operatorname{Hom}_{\mathscr O_Y}(\mathscr Q,\mathscr M).
  \tag{7.7.6.2}\label{III.7.7.6.2}
\]
\end{theorem}

\begin{proof}
En efecto, como $\mathscr M\mapsto
\mathscr F\otimes_{\mathscr O_Y}\mathscr M$ es exacto
\sref[0]{0.6.7.4} y $f_*$ es exacto por la izquierda, el funtor
\[
  \mathscr M\longmapsto
  f_*(\mathscr F\otimes_{\mathscr O_Y}\mathscr M)
\]
es exacto por la izquierda. Basta entonces aplicar la equivalencia de b) y
b$'$) en \sref{III.7.7.5} para $p=1$.
\end{proof}

\begin{corollary}[7.7.7]
\label{III.7.7.7}
Sean $Y$ un preesquema localmente noetheriano, $f:X\to Y$ un morfismo propio,
$\mathscr F$ y $\mathscr F'$ dos $\mathscr O_X$-módulos coherentes y planos
sobre $Y$, y $u:\mathscr F\to\mathscr F'$ un homomorfismo. Consideremos los
dos funtores en el $\mathscr O_Y$-módulo cuasicoherente $\mathscr M$,
\begin{align*}
  \mathscr T(\mathscr M)
  &=
  \operatorname{Ker}
  \bigl(
    f_*(\mathscr F\otimes_{\mathscr O_Y}\mathscr M)
    \longrightarrow
    f_*(\mathscr F'\otimes_{\mathscr O_Y}\mathscr M)
  \bigr),\\
  T(\mathscr M)
  &=
  \Gamma(Y,\mathscr T(\mathscr M))\\
  &=
  \operatorname{Ker}
  \bigl(
    \Gamma(X,\mathscr F\otimes_{\mathscr O_Y}\mathscr M)
    \longrightarrow
    \Gamma(X,\mathscr F'\otimes_{\mathscr O_Y}\mathscr M)
  \bigr).
\end{align*}
Entonces existe un $\mathscr O_Y$-módulo coherente $\mathscr R$ (determinado
salvo isomorfismo único) y existen isomorfismos de funtores
\[
  \mathscr T(\mathscr M)
  \xrightarrow{\sim}
  \mathscr{H}\!om_{\mathscr O_Y}(\mathscr R,\mathscr M)
  \tag{7.7.7.1}\label{III.7.7.7.1}
\]
\[
  T(\mathscr M)
  \xrightarrow{\sim}
  \operatorname{Hom}_{\mathscr O_Y}(\mathscr R,\mathscr M).
  \tag{7.7.7.2}\label{III.7.7.7.2}
\]
\end{corollary}

\begin{proof}
Basta demostrar \eqref{III.7.7.7.2}; en efecto, esto probará
\eqref{III.7.7.7.1} en el caso en que $Y$ es afín, y de ahí se pasará al caso
general razonando como en la demostración de la equivalencia de b) y b$'$) en
\sref{III.7.7.5}, gracias a la unicidad salvo isomorfismo único de un
representante de un funtor representable \sref[0]{0.8.1.8}. De
\sref{III.7.7.6} resulta que existen dos $\mathscr O_Y$-módulos coherentes
$\mathscr Q$, $\mathscr Q'$ que definen isomorfismos funtoriales
\[
  \Gamma(X,\mathscr F\otimes_{\mathscr O_Y}\mathscr M)
  \xrightarrow{\sim}
  \operatorname{Hom}_{\mathscr O_Y}(\mathscr Q,\mathscr M),
  \qquad
  \Gamma(X,\mathscr F'\otimes_{\mathscr O_Y}\mathscr M)
  \xrightarrow{\sim}
  \operatorname{Hom}_{\mathscr O_Y}(\mathscr Q',\mathscr M).
\]
Ahora bien, $u:\mathscr F\to\mathscr F'$ define canónicamente un morfismo de
funtores
\[
  \Gamma(X,\mathscr F\otimes_{\mathscr O_Y}\mathscr M)
  \longrightarrow
  \Gamma(X,\mathscr F'\otimes_{\mathscr O_Y}\mathscr M).
\]

\oldpage[III]{202}
A este último le corresponde un homomorfismo único
$v:\mathscr Q'\to\mathscr Q$ de $\mathscr O_Y$-módulos tal que el diagrama
\[
\xymatrix@C=42pt@R=28pt{
  \Gamma(X,\mathscr F\otimes_{\mathscr O_Y}\mathscr M)
    \ar[r]\ar[d]_-{\sim}
  &
  \Gamma(X,\mathscr F'\otimes_{\mathscr O_Y}\mathscr M)
    \ar[d]^-{\sim}
  \\
  \operatorname{Hom}_{\mathscr O_Y}(\mathscr Q,\mathscr M)
    \ar[r]
  &
  \operatorname{Hom}_{\mathscr O_Y}(\mathscr Q',\mathscr M)
}
\]
sea conmutativo \sref[0]{0.8.1.4}. Como el funtor contravariante
$\mathscr N\mapsto
\operatorname{Hom}_{\mathscr O_Y}(\mathscr N,\mathscr M)$ es exacto por la
izquierda en la categoría de los $\mathscr O_Y$-módulos, basta tomar
$\mathscr R=\operatorname{Coker}(v)$ para obtener el isomorfismo
\eqref{III.7.7.7.2} buscado.
\end{proof}

\begin{corollary}[7.7.8]
\label{III.7.7.8}
Bajo las hipótesis de \sref{III.7.7.6} relativas a $X$, $Y$ y $f$, sean
$\mathscr F$ y $\mathscr G$ dos $\mathscr O_X$-módulos coherentes que verifican
las condiciones siguientes:
\begin{enumerate}
  \item[(i)] $\mathscr F$ es plano sobre $Y$;
  \item[(ii)] $\mathscr G$ es isomorfo al conúcleo de un homomorfismo
    $\mathscr E_1\to\mathscr E_0$ de $\mathscr O_X$-módulos localmente libres
    de tipo finito.
\end{enumerate}
Consideremos los dos funtores en el $\mathscr O_Y$-módulo cuasicoherente
$\mathscr M$,
\begin{align*}
  \mathscr T(\mathscr M)
  &=
  f_*
  \bigl(
    \mathscr{H}\!om_{\mathscr O_X}
      (\mathscr G,\mathscr F\otimes_{\mathscr O_Y}\mathscr M)
  \bigr),\\
  T(\mathscr M)
  &=
  \Gamma(Y,\mathscr T(\mathscr M))
  =
  \operatorname{Hom}_{\mathscr O_X}
    (\mathscr G,\mathscr F\otimes_{\mathscr O_Y}\mathscr M).
\end{align*}
Entonces existe un $\mathscr O_Y$-módulo coherente $\mathscr N$ (determinado
salvo isomorfismo único) y existen isomorfismos de funtores
\[
  \mathscr T(\mathscr M)
  \xrightarrow{\sim}
  \mathscr{H}\!om_{\mathscr O_Y}(\mathscr N,\mathscr M)
  \tag{7.7.8.1}\label{III.7.7.8.1}
\]
\[
  T(\mathscr M)
  \xrightarrow{\sim}
  \operatorname{Hom}_{\mathscr O_Y}(\mathscr N,\mathscr M).
  \tag{7.7.8.2}\label{III.7.7.8.2}
\]
\end{corollary}

\begin{proof}
En virtud del isomorfismo funtorial \sref[0]{0.5.4.2.1}, se tienen, de manera
funtorial en $\mathscr M$, los isomorfismos
\begin{align*}
  \mathscr{H}\!om_{\mathscr O_X}
    (\mathscr E_i,\mathscr F\otimes_{\mathscr O_Y}\mathscr M)
  &\xrightarrow{\sim}
  \check{\mathscr E}_i
    \otimes_{\mathscr O_X}
    (\mathscr F\otimes_{\mathscr O_Y}\mathscr M)\\
  &\xrightarrow{\sim}
  (\check{\mathscr E}_i\otimes_{\mathscr O_X}\mathscr F)
    \otimes_{\mathscr O_Y}\mathscr M\\
  &\xrightarrow{\sim}
  \mathscr{H}\!om_{\mathscr O_X}(\mathscr E_i,\mathscr F)
    \otimes_{\mathscr O_Y}\mathscr M
\end{align*}
para $i=0,1$. Pongamos
$\mathscr F_i=\mathscr{H}\!om_{\mathscr O_X}(\mathscr E_i,\mathscr F)$ para
$i=0,1$; estos son $\mathscr O_X$-módulos coherentes
\sref[0]{0.5.3.5} y planos sobre $Y$ \sref[0]{0.5.4.2}. Si
$v:\mathscr E_1\to\mathscr E_0$ denota el homomorfismo dado, sea
$u=\mathscr{H}\!om(v,1_{\mathscr F}):\mathscr F_0\to\mathscr F_1$. En virtud
de la exactitud por la izquierda del funtor
$\mathscr H\mapsto\mathscr{H}\!om_{\mathscr O_X}
(\mathscr H,\mathscr F\otimes_{\mathscr O_Y}\mathscr M)$, se tienen
isomorfismos funtoriales en $\mathscr M$,
\begin{align*}
  \mathscr{H}\!om_{\mathscr O_X}
    (\mathscr G,\mathscr F\otimes_{\mathscr O_Y}\mathscr M)
  &\xrightarrow{\sim}
  \operatorname{Ker}
  \bigl(
    \mathscr{H}\!om_{\mathscr O_X}
      (\mathscr E_0,\mathscr F\otimes_{\mathscr O_Y}\mathscr M)
    \longrightarrow
    \mathscr{H}\!om_{\mathscr O_X}
      (\mathscr E_1,\mathscr F\otimes_{\mathscr O_Y}\mathscr M)
  \bigr)\\
  &\xrightarrow{\sim}
  \operatorname{Ker}
  \bigl(
    \mathscr F_0\otimes_{\mathscr O_Y}\mathscr M
    \longrightarrow
    \mathscr F_1\otimes_{\mathscr O_Y}\mathscr M
  \bigr).
\end{align*}
Puesto que $f_*$ es exacto por la izquierda, se deduce un isomorfismo funtorial
\[
  f_*
  \bigl(
    \mathscr{H}\!om_{\mathscr O_X}
      (\mathscr G,\mathscr F\otimes_{\mathscr O_Y}\mathscr M)
  \bigr)
  \xrightarrow{\sim}
  \operatorname{Ker}
  \bigl(
    f_*(\mathscr F_0\otimes_{\mathscr O_Y}\mathscr M)
    \longrightarrow
    f_*(\mathscr F_1\otimes_{\mathscr O_Y}\mathscr M)
  \bigr),
\]
y basta entonces aplicar \sref{III.7.7.7}.
\end{proof}

\oldpage[III]{203}
\begin{remarks}[7.7.9]
\label{III.7.7.9}
\textnormal{(i)} En \sref{III.7.7.6}, \sref{III.7.7.7} y
\sref{III.7.7.8}, la formación de los $\mathscr O_Y$-módulos $\mathscr Q$,
$\mathscr R$ y $\mathscr N$ \emph{conmuta con los cambios de base}. Por ejemplo,
conservando las notaciones de \sref{III.7.7.2}, en el caso de
\sref{III.7.7.6} se tiene, para todo $\mathscr O_{Y'}$-módulo cuasicoherente
$\mathscr M'$, el isomorfismo
\[
  f'_*(\mathscr F\otimes_{\mathscr O_Y}\mathscr M')
  \xrightarrow{\sim}
  \mathscr{H}\!om_{\mathscr O_{Y'}}(g^*(\mathscr Q),\mathscr M'),
\]
pues, en virtud de la observación hecha en \sref{III.7.7.2}, todo se reduce a
ver que se tiene
\[
  \operatorname{Hom}_{\mathscr O_Y}(\mathscr Q,g_*(\mathscr M'))
  =
  \operatorname{Hom}_{\mathscr O_{Y'}}(g^*(\mathscr Q),\mathscr M'),
\]
lo cual no es sino \sref[0]{0.4.4.3.1}. Del mismo modo, cuando en
\sref{III.7.7.7} se sustituyen $Y$, $f$, $\mathscr M$, $\mathscr F$ y
$\mathscr F'$ por $Y'$, $f'$, $\mathscr M'$,
$\mathscr F\otimes_{\mathscr O_X}\mathscr O_{X'}$ y
$\mathscr F'\otimes_{\mathscr O_X}\mathscr O_{X'}$, respectivamente, hay que
sustituir $\mathscr R$ por $g^*(\mathscr R)$. Finalmente, en
\sref{III.7.7.8}, cuando se sustituyen $X$, $Y$, $f$, $\mathscr M$ y
$\mathscr F$ por $X'$, $Y'$, $f'$, $\mathscr M'$ y
$\mathscr F\otimes_{\mathscr O_X}\mathscr O_{X'}$, y $\mathscr G$ por
\[
  \mathscr G'=\mathscr G\otimes_{\mathscr O_X}\mathscr O_{X'},
\]
hay que sustituir $\mathscr N$ por $g^*(\mathscr N)$. Esto resulta de que se
tiene todavía una sucesión exacta
\[
  \mathscr E'_1\longrightarrow
  \mathscr E'_0\longrightarrow
  \mathscr G'\longrightarrow0
  \qquad\text{con}\qquad
  \mathscr E'_i=\mathscr E_i\otimes_{\mathscr O_X}\mathscr O_{X'},
\]
y de que
\[
  \check{\mathscr E}'_i
  \otimes_{\mathscr O_{X'}}
  (\mathscr F\otimes_{\mathscr O_Y}\mathscr M')
  =
  \check{\mathscr E}_i
  \otimes_{\mathscr O_X}
  (\mathscr F\otimes_{\mathscr O_Y}\mathscr M')
  \qquad(i=0,1).
\]

\medskip
\textnormal{(ii)} La condición (ii) del enunciado de \sref{III.7.7.8}
relativa a $\mathscr G$ se satisface siempre para cualquier
$\mathscr O_X$-módulo coherente $\mathscr G$ cuando existe un
$\mathscr O_X$-módulo invertible $Y$-amplio, por ejemplo cuando $Y$ es afín y
$f:X\to Y$ es un morfismo proyectivo. Basta observar, teniendo en cuenta
\sref[II]{II.5.5.1}, que existe un $\mathscr O_X$-módulo localmente libre de
tipo finito $\mathscr E_0$ tal que $\mathscr G$ es isomorfo a un cociente de
$\mathscr E_0$ \sref[II]{II.2.7.10}. Como $\mathscr E_0$ y $\mathscr G$ son
coherentes, también lo es el núcleo $\mathscr G_1$ de
$\mathscr E_0\to\mathscr G$; aplicando el mismo resultado a $\mathscr G_1$, se
obtiene una sucesión exacta
\[
  \mathscr E_1\longrightarrow
  \mathscr E_0\longrightarrow
  \mathscr G\longrightarrow0,
\]
donde $\mathscr E_0$ y $\mathscr E_1$ son localmente libres de tipo finito.

\medskip
\textnormal{(iii)} Demostraremos en el capítulo V que, en
\sref{III.7.7.8}, la hipótesis restrictiva (ii) es \emph{superflua}.
\end{remarks}

\begin{proposition}[7.7.10]
\label{III.7.7.10}
\textnormal{(Criterios locales para la propiedad de cambio de base).} Bajo las
condiciones generales de \sref{III.7.7.5}, sean $y$ un punto de $Y$ y $p$ un
entero. Las propiedades siguientes son equivalentes:
\begin{enumerate}
  \item[a)] El funtor $T_p^{\mathscr O_y}$ es exacto por la derecha.

  \item[b)] El homomorfismo canónico
    \[
      T_p^{\mathscr O_y}(\mathscr O_y)
      \longrightarrow
      T_p^{\mathscr O_y}(\kres(y))
    \]
    es sobreyectivo.

  \item[c)] Para todo entero $n$, el homomorfismo canónico
    \[
      T_p^{\mathscr O_y}
      (\mathscr O_y/\mathfrak m_y^{n+1})
      \longrightarrow
      T_p^{\mathscr O_y}(\kres(y))
    \]
    es sobreyectivo.
\end{enumerate}
Además, el conjunto de los $y\in Y$ que verifican estas condiciones es el
mayor abierto $U$ de $Y$ tal que $\mathscr T_p^U$ sea exacto por la derecha.
\end{proposition}

\begin{proof}
Habida cuenta de \sref{III.7.7.4}, la equivalencia de a), b) y c) resulta de
\sref{III.7.4.7} y \sref{III.7.5.2}. El hecho de que el conjunto $U$ donde
$T_p^{\mathscr O_y}$ es exacto por la derecha sea abierto es también
consecuencia de \sref{III.7.4.4}; recíprocamente, si $\mathscr T_p^V$ es exacto
por la derecha, lo mismo ocurre con $T_p^{\mathscr O_y}$ para todo $y\in V$,
por la condición c) de \sref{III.7.7.5} y \sref{III.7.6.1}.
\end{proof}

\begin{corollary}[7.7.11]
\label{III.7.7.11}
Si $\mathscr T_p$ es exacto por la derecha (respectivamente, por la izquierda),
entonces, para todo morfismo $g:Y'\to Y$, $\mathscr T_p^{Y'}$ es exacto por la
derecha (respectivamente, por la izquierda). El recíproco es cierto cuando el
morfismo $g$ es fielmente plano.
\end{corollary}

\oldpage[III]{204}
\begin{proof}
La primera aserción es consecuencia inmediata de \sref{III.7.6.1} y del hecho
de que la cuestión es local sobre $Y$ y $Y'$, por las condiciones c) y b) de
\sref{III.7.7.5}. Para demostrar la segunda aserción basta ver, en virtud de
\sref{III.7.7.10}, que para todo $y\in Y$, $T_p^{\mathscr O_y}$ es exacto por
la derecha (respectivamente, por la izquierda). Pero, por hipótesis, existe
$y'\in Y'$ tal que $g(y')=y$, y $\mathscr O_{y'}$ es un
$\mathscr O_y$-módulo fielmente plano; la conclusión resulta entonces de la
hipótesis y de \sref{III.7.6.1}.
\end{proof}

\begin{remarks}[7.7.12]
\label{III.7.7.12}
\textnormal{(i)} Bajo las hipótesis de \sref{III.7.7.4}, supongamos además que
$\mathscr P_\bullet$ es un complejo \emph{finito}. Resulta entonces de
\sref{III.7.7.1}, puesto que se puede tomar para $\mathfrak U$ un
recubrimiento abierto afín \emph{finito} de $X$, que el bicomplejo
$C^\bullet(\mathfrak U,
\mathscr P^\bullet\otimes_{\mathscr O_Y}\mathscr M)$ es también finito y, más
precisamente, que existe un conjunto finito $E$ de pares, \emph{independiente
de $\mathscr M$}, tal que
\[
  C^h(\mathfrak U,\mathscr P^k\otimes_{\mathscr O_Y}\mathscr M)=0
  \qquad\text{para todos los pares }(h,k)\notin E.
\]
Se concluye que existe $i_1$ tal que, para $i\geq i_1$, se tiene
$\mathscr T_i(\mathscr M)=0$ para \emph{todo} $\mathscr O_Y$-módulo
cuasicoherente $\mathscr M$. En particular, $\mathscr T_i$ es trivialmente un
funtor exacto en $\mathscr M$ para estos valores de $i$ y, por consiguiente,
por \sref{III.7.4.1}, $Z'_i(L_\bullet)$ es un $A$-módulo \emph{plano} de tipo
finito (y, por tanto, \emph{proyectivo} de tipo finito, puesto que $A$ es
noetheriano) para estos valores de $i$.

Consideremos entonces el complejo $L'_\bullet$ de $A$-módulos tal que
\[
  L'_i=L_i\quad(i<i_1),
  \qquad
  L'_{i_1}=Z'_{i_1}(L_\bullet),
  \qquad
  L'_i=0\quad(i>i_1),
\]
y sea $\mathscr L'_\bullet=\widetilde{L'_\bullet}$. Es claro que
\[
  \mathscr H_i(\mathscr L'_\bullet\otimes_{\mathscr O_Y}\mathscr M)
  =
  \mathscr H_i(\mathscr L_\bullet\otimes_{\mathscr O_Y}\mathscr M)
\]
para $i<i_1-1$ y también para $i\geq i_1$ (siendo entonces nulos ambos
miembros). Finalmente, como
\[
  \operatorname{Im}(Z'_{i_1}\otimes_A M)
  =
  \operatorname{Im}(L_{i_1}\otimes_A M)
\]
por definición, se tiene asimismo
\[
  \mathscr H_i(\mathscr L'_\bullet\otimes_{\mathscr O_Y}\mathscr M)
  =
  \mathscr H_i(\mathscr L_\bullet\otimes_{\mathscr O_Y}\mathscr M)
\]
para $i=i_1-1$. Se ve, pues, que se puede suponer en \sref{III.7.7.4} que
$\mathscr L_\bullet$ es también un complejo \emph{finito}, a condición de
exigir solamente que los $\mathscr L_i$ sean $\mathscr O_Y$-módulos
\emph{localmente libres}, asociados a $A$-módulos proyectivos de tipo finito.

Este razonamiento se aplica, en particular, al caso en que
$\mathscr P_\bullet$ se reduce a \emph{un solo} término $\mathscr F\ne0$, de
grado $0$ (en cuyo caso
$\mathscr T_n(\mathscr M)=\mathrm R^{-n}f_*
(\mathscr F\otimes_{\mathscr O_Y}\mathscr M)$). Se puede entonces suponer que
los $\mathscr L_i$ son nulos para $i>0$; en este caso se utilizarán
preferentemente las notaciones cohomológicas, escribiendo $\mathscr T^{-p}$ en
lugar de $\mathscr T_p$.

\medskip
\textnormal{(ii)} Cuando en el enunciado de \sref{III.7.7.5} ya no se supone
que los $\mathscr P_i$ sean planos sobre $Y$, las conclusiones siguen siendo
válidas a condición de poner esta vez
\[
  \mathscr T_p(\mathscr M)
  =
  \operatorname{Tor}^Y_p
    (f,1_Y;\mathscr P_\bullet,\mathscr M).
  \tag{7.7.12.1}\label{III.7.7.12.1}
\]
En efecto,
$\operatorname{Tor}^Y_n(f,1_Y;\mathscr P_\bullet,\mathscr O_Y)$ es entonces un
$\mathscr O_Y$-módulo \emph{coherente} en virtud de \sref{III.6.7.9}. La
demostración de \sref{III.6.10.5} se aplica sin cambio, teniendo en cuenta
\sref{III.6.10.1}, y muestra también que, cuando
$Y=\operatorname{Spec}(A)$ es afín, se tiene
\[
  \mathscr T_p(\mathscr M)
  =
  \mathscr H_p(\mathscr L_\bullet\otimes_{\mathscr O_Y}\mathscr M),
  \qquad
  \mathscr L_\bullet=\widetilde{L_\bullet},
\]
donde $L_\bullet$ es un complejo de $A$-módulos libres de tipo finito; esto
prueba la aserción.
\end{remarks}
